\scanpage{32}
\mspage{12}
following properties (characterizing a \uline{strict} Quillen triple)

4) \deleted{A any cobase change} \inserted{$W$ stable by base} change by any $\deleted{\unsure{g}}f\in F$, and by cobase change by any $\deleted{\unsure{f}}g\in C$.

\deleted{\unsure{5)}} \uline{NB} From my point of view, $W$ is given beforehand in a rather tangible way, by \unsure{cohomological|homological} criteria, which turns out \unsure{very|easy} handy. In Kan's--Quillen's point of view, $W$ was a little more hidden maybe in terms of (say) ss structures, they get to it via \deleted{\unsure{a}} the pair $(C,TF)$ and $F\supset TF$ the Kan fibrations, which is a somewhat ad-hoc construction in terms of a certain \deleted{\unsure{g}} a-priori generating set $(TC)_0$ --- corresponding roughly to expressing the horn lifting property. What remains to be understood, in the context say of general test categories, is the relationship between this and the direct description $TC=C\cap W$, \deleted{\unsure{what}} i.e why this is just the \uline{saturation} of $(TC)_0$. Maybe I should have a look at Gabriel-Zisman and their treatment of ``anodyne extensions''.
% archive sheet number: 31
